\documentclass[14pt]{article} 
\usepackage[utf8]{inputenc}
\usepackage{epsfig}
\usepackage[ngerman]{babel}
\usepackage{mathtools,extarrows}
\usepackage[thmmarks]{ntheorem}
\usepackage{amsmath,theorem}
\usepackage{amsfonts}
\usepackage{amssymb}
\usepackage{txfonts}
\usepackage{marvosym}
\usepackage{MnSymbol}
\usepackage{mathrsfs}
\usepackage{parallel}
\usepackage{tikz}
\usepackage{geometry}
\usepackage{titletoc}
\usepackage{url}
\usepackage[pdftex,bookmarks=true,bookmarksnumbered=true,bookmarksopen=true,colorlinks=true,
		filecolor=black,linkcolor=blue,urlcolor=blue,plainpages=false,pdfpagelabels,
		citecolor=black,setpagesize=false]{hyperref}


\makeatletter
\renewcommand\l@section{\@dottedtocline{1}{0em}{2.5em}}
\renewcommand\l@subsection{\@dottedtocline{2}{2.5em}{2em}}
\makeatother

\titlecontents{chapter}[1.5em]{\addvspace{1em}\bfseries}
	{\contentslabel{1.5em}}{\hspace^{-1.5em}}
	{\hfill\contentspage}
\dottedcontents{section}[3.8em]{}{1.3em}{1pc}
\dottedcontents{subsection}[6.0em]{}{2.2em}{1pc}
\dottedcontents{subsubsection}[9.2em]{}{3.1em}{1pc}

%%%%theoremenvironments
\theoremstyle{plain} 

\theoremheaderfont{\bfseries}

\theorembodyfont{\it}
\newtheorem{theorem}{Theorem}[section]
\newtheorem{lem}[theorem]{Lemma}
\newtheorem{cor}[theorem]{Corollar}
\newtheorem{prop}[theorem]{Proposition}

\theorembodyfont{\rm}
\newtheorem{remark}[theorem]{Bemerkung}
\newtheorem{defi}[theorem]{Definition}
\newtheorem{rul}[theorem]{Regel}

\theoremstyle{nonumberplain}

\theoremheaderfont{%
    \normalfont\scshape}
\theorembodyfont{\normalfont}
\theoremsymbol{\square}
\newtheorem{proof}{Beweis:}

\topmargin=-2mm \textwidth=155mm \textheight=230mm
\oddsidemargin=0mm \evensidemargin=0mm
\thispagestyle{empty}
\newcommand{\mb}{\mathbb}
\newcommand{\mc}{\mathcal}
\newcommand{\mf}{\mathfrak}
\newcommand{\mr}{\mathrm}
\newcommand{\ms}{\mathscr}
\newcommand{\K}{\mb{K}}
\newcommand{\C}{\mb{C}}
\newcommand{\R}{\mb{R}}
\newcommand{\Q}{\mb{Q}}
\newcommand{\Z}{\mb{Z}}
\newcommand{\N}{\mb{N}}
\newcommand{\ra}{\rightarrow}
\newcommand{\LRa}{\Leftrightarrow}
\newcommand{\Ra}{\Rightarrow}
\newcommand{\La}{\Leftarrow}
\newcommand{\eps}{\varepsilon}
\newcommand{\vt}{\vartheta}
\newcommand{\vp}{\varphi}
\newcommand{\qq}{\hspace{.25cm}}
\newcommand{\ol}{\overline}
\newcommand{\ul}{\underline}
\renewcommand{\d}{ \, \mr{d}}
\newcommand{\Res}{\mr{Res}}
\newcommand{\Const}{\mr{Const} \, }
\newcommand{\Id}{\mr{Id}}
\geometry{hmargin=3cm,top=2cm,bottom=4cm}

\newcommand{\makelicense}{Dieser Mitschrieb steht unter der freien \texttt{CC-BY-SA-DE 3.0} 
Lizenz.
\begin{center}\includegraphics[height=31pt,width=88pt]{license.pdf}\end{center}
Für weitere Informationen besuchen Sie
\begin{center}\url{http://creativecommons.org/licenses/by-sa/3.0/deed.de}\end{center}}


\input xy
\xyoption{all}

\title{\sc{Nichtlineare Funktionalanalysis}}
\author{Luis Felipe Müller}
%\date{Sommersemester 2011}

\begin{document}


\maketitle

\hrule \hspace*{1cm}\\[3mm]
{\sc Kompletter Mitschrieb zur gleichnamigen Vorlesung bei Herrn Dondl\\
\hspace*{1cm}\hfill (Sommersemester 2011, Uni Heidelberg)}\\
\hrule
\vspace{1cm}
\noindent\makelicense

%\vfill
%\listoffigures
%\vfill
%\listoftables

\newpage

\pagenumbering{roman}

\section*{Organisatorisches}

\subsection*{Termine}
\begin{itemize}
    \item Vorlesung: Mo./Mi. 9-11ct. -104 Ang. Math.
    \item Übungsaufgaben: Mi. - Mi. vor der Vorlesung. Kasten in der Angew.Math
    \item Übungsgruppe: Fr. 16-18 Angew. Math. -101
    \item Website zur Vorlesung: \url{http://dondl.org/wiki/Sommersemester_11}
    \item Literatur: \begin{enumerate}
        \item Růžička M: Nichtlineare Funktionalanalysis, Eine Einführung
        \item Aubin-Ekeland: Applied nonlinear Analysis
        \item Deimling: Nonlinear Functional Analysis
        \item Schwartz: Nonlinear Functional Analysis
        \item Zeidler: Nonlinear Functional Analysis and its applications 
        \end{enumerate}
    \item Prüfungen: Zulassung mit etwa 50\% der Übungsaufgaben-Punkte. Prüfung ist mündlich, 
            beispielsweise am 29. Juli (Fr.) 
    \item Dozent: Patrick Dondl, Sprechstunde Mo./Mi. 11-12 in Raum 130 (Angew. Math.)
    \item Tutor: Julian Scheuer
\end{itemize}

\newpage

\tableofcontents

\newpage

\pagenumbering{arabic}

\chapter{Einleitung}

\section{Thema der Vorlesung}

In der linearen Funktionalanalysis haben wir eine Vielzahl von Methoden kennengelernt um
Ergebnisse aus der endlichdimensionalen linearen Algebra auf den unendlichdimensionalen Fall 
zu verallgemeinern. Ein Hauptaufgabe war dabei, die Lösbarkeit von Gleichungen der Form
\[
    Ax=y
\]
für lineare Operatoren $A$ auf $\infty$-dimensionalen Banchräumen zu zeigen.
\begin{description}
    \item{Ein Beispiel:}
    
    Sei $\Omega\subset\R^3$, gefüllt mit einer inkompressiblen, riskosen Flüssigkeit.
    $v_j(x)$ sei die Geschwindigkeit der Flüssigkeit an der Stelle $x\in \Omega$. 
    $p(x)$ ist der Druck an der Stelle $x$.
    \begin{description}
    \item{Randbedingungen:}
    $v_j(x)=0 \qq x\in \partial\Omega$
    \item{Inkompressibilität:}
    $\partial_jv_j(x)=0 \qq x\in \Omega$
    \item{Bewegungsungleichung:}
    Wir betrachten die Kräfte, die auf einen kleinen Würfel, eingeschlossen durch
    $(x_1,x_2,x_3)$, $(x_1+\Delta x_1, x_2+\Delta x_2, x_3+\Delta x_3)$. Druck auf eine Oberfläche
    des Würfels mit Normale $x$, ist $f_j^p=p\cdot \Delta x_2 \Delta x_3 \cdot \delta_{ij} $ auf die
    gegenüberliegende Seite wirkt
    \[
        -(p+\partial _i p \Delta x_i) \Delta x_2 \Delta x_3 \delta_{ij}
    \]
    \end{description}
Zusammen ergibt sich 
\[
    f=(\partial_j p)\Delta V
\]

Kraft durch Viskosität auf eine Oberfläche mit Normale $x_1$ ist 
\[
    f_j^{V,x_i}= -2\eta \Delta x_2\Delta x_3 \partial_1v_j
\]
mit einer Konstante $\eta$. Der gleiche Trick wie oben ergibt für die gegenüberliegende
Oberfläche
\[
    \eta \Delta x_2 \Delta x_3 \partial (v_j+\partial v_j\Delta x_1)
\]
Zusammen ergibt sich
\begin{align*}
    f_j^V&= \eta\Delta V \cdot \partial_i\partial_iv_j\\
    \text{Newton:} \qq \rho \Delta V \frac{\d}{\d t} v_j (t,x(t))
        &= \eta \Delta V  \partial _i\partial_iv_j- \Delta V (\partial _j p) + \Delta 
        V\kappa_j \qq \text{mit einer externen Konstante $\kappa_j$, bspw. Gravitation}
\end{align*}
Teilen durch $\Delta V$ und die Kettenregel ergibt
\begin{align*}
    \rho \partial_t v_j&= \eta \partial_i\partial_i v_j - \rho (v_i\partial_i) - \partial_j p + K_j 
        \qq \text{(Newton)}\\
    d_jv_j&=0
\end{align*}
Naurier-Stokes. Frage: Existiert eine eindeutige Lösung zur sationären Naurier-Stokes-Gleichung:
\begin{align}\label{1}
    \eta \partial_i\partial_iv_j-(v_i\partial_i)v_j+\partial_jp+K_j=0
\end{align}
\end{description}
Wir können die Gleichung etwas umschreiben. Sei $H$ ein Hilbertraum
\[
    H:=\ol{\left\{ u\in C^\infty_c (\Omega,\R^3), \text{ so dass } \partial_jv_j=0 \right\}}
        ^{W^{1,2}(\Omega, \R^3)}
\]
Ein Skalarprodukt auf $H$ ist gegeben durch
\begin{align*}
    (u,v)_H:= \int_{\Omega} \nabla u\cdot \nabla v \d x = \int_\Omega \sum_{i,j=1}^3 
    (\partial_iu_j)(\partial_iv_j)
\end{align*}
$\Omega$ beschränkt $\Ra$ $(\cdot, \cdot)_H$ ist äquivalent zum üblichen Skalarprodukt 
(mittels Poincaré).
Wir multiplizieren (\ref{1}) mit $\omega \in H $, integrieren und erhalten
\begin{align*}
    \int_\Omega (\eta \partial_i\partial_i v_j- (v,\partial_i) v_j + K_j) \cdot \omega_j 
    &= \int_\Omega (\partial _jp)\omega_j=0 \text{, da $\omega$ divergenzfrei.}\\
    (\ref{1}) \Ra \eta (v,\omega)_H-a(v,v,w)-\int_\Omega K\omega&=0
\end{align*}
Ebenso für $a$:
\begin{align*}
    a(u,v,w):= (\underbrace{B(u,v)}_{\text{bilinear.}},w)_H
\end{align*}
Also
\begin{align*}
    (\ref{1}) \Ra \, (\eta v - B(u,v)) - \tilde K , \omega)_H=0 \qq \forall \omega \in H
\end{align*}
somit
\[
    \eta v - B(v,v) = \tilde K
\]
Das ist eine Gleichung der Form
\begin{align}\label{2}
    Fv=\tilde K,\qq\text{mit $F$ einem Nichtlinearen Operator}
\end{align}
Im ersten Teil der Vorlesung beschäftigen wir uns mit der eindeutigen Lösbarkeit von
Gleichungen der Form
\[
    Fx=y, \qq F:X\ra Y, \qq X,Y\qq \text{Banachräume}
\]
und zum Abschluß zeigen wir mit Hilfe des \textit{Schauder'schen Fixpunktsatzes} die Existenz
und finden eine Lösung von (\ref{2}), also der schwachen Form der stationären 
Navier-Stokes-Gleichung.

Im zweiten Teil der Vorlesung beschäftigen wir uns mit Variationsrechnung (d.h. dem Finden
von Minimierern nichtlinearer Funktionalen)
\[
    W:X\ra \R \qq\text{mit} \qq X \qq \text{ein Banachraum}
\]
Finde 
\[
    x_0\in X: \, W(x_0)= \inf_{y\in X}W(y)
\]
Insbesonder treffen wir dort auf Probleme in der Elastizitätstheorie.

\section*{Aufbau der Vorlesung}
\begin{itemize}
    \item Abbildungsgrad $\ra$ Existenz von Lösungen von $Fx=y$
    \item Monotone Operatoren $\ra$ Eindeutigkeit von Lösungen von $Fx=y$; 
    	zeitabhängige Probleme.
    \item Variationsrechnung $\ra$ $\inf_{y\in X} W(y)$
\end{itemize}

\subsection{Vorarbeiten}

\subsubsection{Ableitung in Banachräume und implizite Funktionen}

Es seien $X$ und $Y$ Banachräume, $\Omega\subset X$ offen, $F:\Omega \ra Y$, $x_0\in \Omega$

\begin{defi}[Gâteaux-Ableitung]
    Die Gâteaux-Ableitung $\d F(x_0,\psi)$ des Operators $F$ im Punkt $x_0$ in Richtung $\psi\in X$ 
    ist gegeben durch 
    \[
        \d F (x_0,\psi)=\lim_{s\ra 0} \frac{F(x_0+s\cdot \psi)-F(x_0)}{s}= \left. \frac{\d}{\d s} 
        F(x_0+s\psi)\right|_{s=0}
    \]
    falls der Limes existiert. Der Operator $F$ heißt in diesem Fall in $x_0$ Richtung $\psi$ 
    Gâteaux-differenzierbar.
\end{defi}
\begin{defi}[Fréchet-Ableitung]
    Der Operator $F$ heißt Fréchet-differenzierbar in $x_0\in \Omega$, falls ein beschränkter,
    linearer Operator 
    \[
        F'(x_0):X\ra Y 
    \]
    existiert, so dass 
    \begin{align}\label{3}
        \lim_{\|h\|\ra 0} \frac{\| F(x_0+h) - F(x_0) -F'(x_0)\cdot h \|  }{\|h\|} = 0
    \end{align}
    $F'(x_0)$ heißt dann Fréchet-Ableitung von $F$ in $x_0$.
\end{defi}

\begin{theorem}
    \begin{description}
        \item{i)} $F'(x_0)$ ist durch (\ref{3}) eindeutig bestimmt.
        \item{ii)} Falls $F$ stetig ist in $x_0$, so ist jeder lineare Operator, der (\ref{3}) erfüllt,
                ebenfalls stetig.
        \item{iii)} Ist $L:X\ra Y$ linear, so gilt
                \[
                    L'(x)=L \qq \forall x\in X
                \]
    \end{description}
\end{theorem}

\begin{proof}

\begin{description}

\item{i)} 
Es gelte (\ref{3}) auch für L. Dann haben wir
\begin{align*}
    \|Lh-F'(x_0)h\| \leq \eps \|h\| \qq \text{falls} \qq \|h\| < \delta = \delta(\eps, x_0)
\end{align*}
Für beliebiges $h$ folgt aber
\begin{align*}
    \|(L-F'(x_0)) (\delta\|h\|^{-1}\cdot h)\|&\leq \delta\eps\\
    \Ra \, \|(L-F'(x_0))h  \|&\leq \eps \|h\| \qq \forall \, h\in X, \, \forall \, \eps>0 \\
    \Ra \, \|L-F'(x_0)\|_{\ms L(X,Y)} &= 0
\end{align*}

\item{ii)}
(\ref{3}) wird umgeformt zu
\[
    \|F'(x_0) h \| \leq \eps \|h\| + \|F(x_0+h)-F(x_0)\|
\]
Mit $h\ra 0$ folgt die Stetigkeit (für $\|h\|\leq \delta$) von $F'(x_0)$ an der Stelle 0. 
Wegen Linearität von $F'(x_0)$ ist $F'(x_0)$ somit stetig. 

\item{iii)} (\ref{3}) gilt offensichtlich für $L'(x_0)=L$, mit i) folgt Eindeutigkeit.
\end{description}
\[ \]
\end{proof}


\begin{prop}
    Jeder Fréchet-differenzierbare Operator $F$ ist Gâteaux-differenziebar $\forall \, \psi\in X$ 
    und es gilt
    \[
        F'(x_0)\psi=\d F(x_0,\psi)
    \]
\end{prop}
\begin{proof}
    Übungsaufgabe \[\]
\end{proof}

\begin{defi}
    $F$ heißt (Fréchet-)differenzierbar auf $\Omega$, falls $\forall\, x\in X$ ein $F'(x)$ existiert,
    sodass $F'(x)$ stetig ist  und (\ref{3}) erfüllt. $F$ heißt stetig (Fréchet-)differenzierbar in
    $\Omega$, falls die Abbildung
    \[
        F': \Omega\ra \ms L(X,Y)
    \]
    stetig ist.
\end{defi}

\begin{prop}
    Existiert die Gâteaux-Ableitung $\d F(x,\psi) \, \forall x\in \Omega$, und ist sie linear und
    stetig in $\psi$ $\forall\, x\in \Omega$, so ist $F$ Fréchet-differenziebar auf $\Omega$ und
    es gilt 
    \[
        F'(x)\psi=\d F(x,\psi)
    \]
\end{prop}

\begin{proof} 
    Übungsaufgabe \[\]
\end{proof}

\begin{defi}
    Sei $F$ auf $\Omega$ stetig differenziebar, $x_0\in \Omega$. Falls ein stetiger linearer Operator
    \[
        F''(x_0):X\ra \ms L(X,Y)
    \]
    existiert mit
    \[
        \lim_{\|h\|\ra 0}\frac{\|F'(x_0+h)-F'(x_0)-F''(x_0)h\|_{\ms L(X,Y)}}{\|h\|}=0
    \]
    dann heißt $F$ in $x_0$ zweimal (Fréchet-)differenzierbar und $F''(x_0)$ heißt zweite Ableitung
    von $F$ in $x_0$. Höhere Ableitungen entsprechend.
\end{defi}

\begin{remark}
    Es gilt die Kettenregel:
    Seien $X$, $Y$, $Z$ Banachräume, $\Omega_X\subset X$ offen, $x_0\in \Omega_X$,
    \[
        F:\Omega_X\ra Y, \qq F(x_0)=y_0\in \Omega_Y\subset Y \qq \text{offen}
    \]
    \[
        G:\Omega_Y\ra Z
    \]
    Falls $F'(x_0)$ und $G'(y_0)$ existiert, so ist 
    \[
    (G\circ F)'(x_0)=G'(y_0)\circ F'(x_0)
    \]
\end{remark}

\begin{defi}[Partielle Ableitung]
    Seien $X,$ $Y$, $Z$ Banachräume, $\Omega_X\subset X$ offen, $x_0\in \Omega_X$, $\Omega_Y\subset Y$ 
    offen, $y_0\in \Omega_Y$. Der Operator
    \[
        F: \Omega_X \times \Omega_Y \ra Z 
    \]
    heißt partiell in $(x_0,y_0)$ nach dem zweiten Argument (nach $y$) differenzierbar, falls die 
    Ableitung 
    \[
        F(x_0, \cdot) : \Omega_Y \ra Z 
    \]
    differenziebar ist. Wir nennen den linearen Operator $F_Y(x_0,y_0): Y\ra Z$, der
    \[
        \lim_{\|h\|\ra 0} \frac{\| F(x_0,y_0+h) - F(x_0,y_0) - F_Y( x_0,y_0)h \|}{\|h\|} = 0 
    \]
    erfüllt, die partielle Ableitung von $F$ in $(x_0,y_0)$ nach dem zweiten Argument.
\end{defi}

\begin{prop}
    Seien $X,Y$ Banachräume, $\Omega \subset X $ offen und konvex mit $x_0, x_1\in \Omega$. 
    $F: \Omega \ra Y$ sei stetig Fréchet-differenzierbar auf $\Omega$. Dann gilt
    \[
        F(x_1)-F(x_0)=\int_{0}^1 F'(x_0+t(x_1-x_0))(x_1-x_0)\d t
    \]
    Das Integral ist als Limes der entsprechenden Riemannsumme zu verstehen und dieser existiert.
\end{prop}

\begin{proof}
    Übungsaufgabe \[\]
\end{proof}

Ähnlich dem endlichdimensionalen Fall geben uns die Ableitungen im Banachraum
hinreichende Bedingungen um Operatoren implizit zu definieren. Die Fragestellung ist die folgende:
Seien $X,\, Y,\, Z$ Banachräume, $U$ eine Umgebung von $x_0\in X$. $V$ eine Umgebung von $y_0\in Y$.
Wir suchen zu $F:U\times V \ra Z$ einen Operator
\[
    T:U_0\subset U \ra V, 
\]
sodass gilt
\[
    F(x,Tx)=F(x_0,y_0) \qq \forall \, x\in U_0.
\]
Durch eine einfache Verschiebung ist es ausreichend, den Fall
\[
    F(x_0,y_0)=0
\]
zu untersuchen.

\begin{prop}
    Sei $X$ ein Banachraum, $\Id: X\ra X$, $x\mapsto x$ die Identität auf $X$. Es sei
    \[
        R:B_r(0)\subset X \ra X, 
    \]
    eine $k$-Kontraktion mit $k<1$, d.h. $\|R(x)-R(y)\|\leq k \|x-y\|$, und es gelte
    \[
        \|R(0)\| < r(1-k)
    \]
    Dann existiert genau ein $x\in B_r(0)$ mit
    \[
        (\Id + R)x = 0
    \]
\end{prop}

\begin{proof}
    Sei $S=-R$, wir suchen also einen Fixpunkt von $S$.
    \begin{description}
        \item{1. Eindeutigkeit:} 
        Seien $Sx=x$ und $Sx'=x'$, damit gilt
        \[
            \|x-x'\|=\|Sx-Sx'\| \leq k \|x-x'\|.
        \]
        Mit $k<1$ folgt $x=x'$.
        
        \item{2. Existenz:}
        Sei $x\in B_r(0)$, es gilt
        \[
            \|Sx\|\leq \|Sx-S(0)\|+\|S(0)\|\leq k\|x\|+\|S(0)\|< kr + r(1-k) = r
        \]
        Sei $x_p=Sx_{p-1}$, $x_0=0$. Es gilt (wie auch im Banach'schen Fixpunktsatz, siehe ÜB 1), dass
        \[
            \|x_{n+p}-x_n\| \leq k^n(1-k)^{-1}\|x\|,
        \]
        damit ist $(x_p)_{p\in \N}$ eine Cauchy-Folge und konvergiert gegen $x\in X$.
        Wir haben weiter, dass
        \[
            \|x\| \leq \underbrace{\|x-x_{p+1}\|}_{\ra 0} + \|x_{p+1}\| \qq \text{mit} 
            \qq \|x_{p+1}\|< (1-k) \|S(0)\|=r\\
            \Ra \, \|x\|<r
        \]
        Wegen $x_{p+1}=Sx_p$ gilt dass $x=Sx$, somit ist $x$ der gesuchte Fixpunkt. 
    \end{description}
    \[  \]
\end{proof}

\end{document}
